\chapter{Pierwsze kroki w teorii reprezentacji}

% POTEM: Diagramy younga

Reprezentacja grupy to próba zrozumienia jej przez pryzmat przestrzeni wektorowej.

\section{Podstawowe definicje}

\begin{definition}{reprezentacja liniowa}{}
  Niech $G$ będzie grupą, a $V$ przestrzenią wektorową nad ciałem $K$ (domyślnie $K=\C$). \buff{Liniową reprezentacją} grupy $G$ na $V$ nazywamy działanie grupy $G$ na $V$ przez liniowe automorfizmy. Innymi słowy, reprezentacja to liniowy homomorfizm
  $$\rho:G\to GL(V).$$
\end{definition}

Jeśli $\dim(V)=n$ to przez $GL(V)$ rozumiemy macierze $n\times n$ o wyrazach z $K$ i niezerowym wyznaczniku.

\begin{example}[m]
  \item Rozważmy odwzorowanie 
    $$\rho: \Z_n\to GL(2, \R)$$
    zadane przez 
    $$\rho(g)=\begin{bmatrix}\cos(\frac{2\pi}{g})&-\sin(\frac{2\pi}{g})\\ \sin(\frac{2\pi}{g})&\cos(\frac{2\pi}{g})\end{bmatrix}.$$
    Od razu poznajemy, że macierz $\rho(g)$ to macierz obrotu o kąt $\frac{2\pi}{g}$. Nad $\C$ możemy zapisać $\rho(g)=e^{\frac{2\pi i}{g}}$.
  \item Popatrzmy teraz na dwie reprezentacje znanej nam już grupy permutacji $S_3$. 

    Pierwsza reprezentacja to działanie na $\R^3$ przez permutowanie współrzędnych. Taka reprezentacja ma dwuwymiarową podprzestrzeń niezmienniczą: $V=\{(x_1,x_2, x_3)\;:\;x_1+x_2+x_3=0\}$.

    Zanim zdefiniujemy drugą reprezentację, przygotujemy pasującą nam przestrzeń liniową. Niech $V$ będzie rzeczywistą przestrzenią wektorową z bazą $\{e_{\sigma_1}, e_{\sigma_2}\}$ i symetryczną formą dwuliniową
    $$B(e_{\sigma_i}, e_{\sigma_j})=\begin{cases}1 & i=j\\\cos(\frac{\pi}3)&i\neq j\end{cases}.$$
    Definiujemy reprezentację (nazywaną \emph{reprezentacją Titsa}) na generatorach $S_3$ jako $\rho(\sigma_i)(v)=v-2B(v, e_{\sigma_i})e_{\sigma_i}$. Sprawdzenie poprawności oraz czy autorka nie pomyliła znaków zostaje pozostawione jako ćwiczenie dla czytelnika.
\end{example}

\begin{definition}{równoważność reprezentacji}{}
  Niech $\rho:G\to GL(V)$ i $\rho':G\to GL(V')$ będą dwoma reprezentacjami grupy $G$. Powiemy, że $\rho$ i $\rho'$ są \buff{równoważne} (lub izomorficzne), \buff{$\rho\approx\rho'$}, jeśli istnieje izomorfizm $\phi: V\to V'$ taki, że diagram
  $$\begin{tikzcd}[column sep=large]
    V\arrow[r, "\phi"]\arrow[d, "\rho(g)" left] & V'\arrow[d, "\rho'(g)"]\\ 
    V\arrow[r, "\phi" below] & V'
  \end{tikzcd}$$
  komutuje dla wszystkich $g$.
\end{definition}

\begin{example}
  $G$ - grp skończona, $l^2(G)\ni f$, $(\rho(g)f)(h)=f(hg)$

  wszystkie parami nieizomorficzne nieprzywiedlnie reprezentacje $\{(\rho_1, V_1),...,(\rho_k, V_k)\}$ $G$
  $\pi: V_1^{\dim V_1}\oplus...\oplus V_k^{\dim V_k}$

  wtedy $\rho\simeq \pi$


  Dowód zaczyna od pokazania, że $k[G]\cong \oplus_{n} M_n(C)$

  % \color{red} reprezentacje Titsa na różnych wyborach generatorów?
\end{example}

Dla ustalonej grupy $G$ możemy rozważać kategorię jej reprezentacji liniowych. Obiektami będą wtedy pary $(V, \rho)$ (zwykle jednak będziemy odwzorowanie i przestrzeń utożsamiać). Jako morfizmy między reprezentacjami zdefiniujemy \buff{przeplatacze}, czyli liniowe odwzorowania spełniające diagram wyżej, ale niekoniecznie będące izomorfizmami. Dla $\rho:G\to V$, $\rho':G\to V'$ zbiór przeplataczy między nimi będziemy czasem (zależnie od wybryków autorki) oznaczać $Hom_G(\rho, \rho')$, a czasami $Hom_G(V,V')$.

\subsection{Reprezentacje nieprzywiedlnie}

\begin{definition}{reprezentacja nieprzywiedlnia}{}
  Niech $\rho:G\to GL(V)$ będzie reprezentacją liniową. Powiemy, że $\rho$ jest \buff{nieprzywiedlnia} (ang. irreducible), jeśli $V$ nie zawiera żadnej niezmienniczej na działanie $G$ przez $\rho$ podprzestrzeni różnej od $0$ i całego $V$.
\end{definition}

Dla przypomnienia: podprzestrzeń $W\subseteq V$ jest niezmiennicza na działanie $G$, jeśli dla każdego $g\in G$ zachodzi $gW\subseteq W$.

\begin{example}
  Reprezentacja $S_3$ przez macierze permutacji nie jest nieprzywiedlnia na $\R^3$, ale jeśli ograniczymy tę reprezentację do podprzestrzeni $V_0=\{(x_1,x_2,x_3)\;:\;\sum x_i=0\}$ to staje się ona nieprzywiedlnia.
\end{example}

Załóżmy, że grupa $G$ działa na przestrzeni wektorowej $V$ z iloczynem skalarnym $\langle\cdot,\cdot\rangle$ w sposób zachowujący ten iloczyn. To znaczy, dla każdego $g\in G$ oraz $v,w\in V$ zachodzi $\langle v,w\rangle=\langle gv,gw\rangle$. W takiej sytuacji reprezentacja grupy $G$ dana przez to działanie jest przywiedlnie wtedy i tylko wtedy gdy istnieje rozkład $V=W\oplus W^\perp$ taki, że zarówno $W$ jak i $W^\perp$ są niezmiennicze na $G$.

Dla grup skończonych działających na przestrzeni $(V, \langle\cdot,\cdot\rangle)$ zawsze jesteśmy w stanie stworzyć iloczyn skalarny niezmienniczy na jej działanie: 
$$\langle v,w\rangle_G=\frac{1}{|G|}\sum_{g\in G}\langle gv, gw\rangle.$$

W bardzo podobny sposób jesteśmy w stanie stworzyć niezmienniczy iloczyn skalarny na grupie zwartej (wystarczy wziąć do ręki miarę Haara $G$ zamiast $|G|$ i wycałkować). Dla grup lokalnie zwartych jest to już nieco problematyczne, gdyż miara Haara całej grupy wcale nie musi być skończona, wystarczy popatrzeć na $\R$.

\begin{theorem}{}{}
  Jeśli $G$ jest grupą skończoną (lub zwartą), to każda jej reprezentacja $\rho:G\to GL(V)$ jest sumą prostą reprezentacji nieprzywiedlnich.
\end{theorem}

\begin{proof}
  Udowodnimy twierdzenie wyżej przez indukcję względem $\dim(V)$.

  Dla $\dim(V)=0$ wszystko jest jasne. Dla $\dim(V)\geq 1$ sprawdzamy, czy istnieje podprzestrzeń $W\subseteq V$ niezmiennicza na $\rho$. Jeśli nie istnieje, to $\rho$ jest już nieprzywiedlnia. W przeciwnym przypadku na mocy uwagi wyżej mamy rozkład $V=W\oplus W^\perp$. Ponieważ $W$ jak i $W^\perp$ są niezmiennicze na działanie $G$, możemy obciąć $\rho$ do tych podprzestrzeni tak, że $\rho=\rho\restriction{W}\oplus \rho_{\upharpoonright W^\perp}$
  Ponieważ $\dim(W), \dim(W^\perp)<\dim(V)$ to znajdujemy rozkład $\rho_{\upharpoonright W}$ oraz $\rho_{\upharpoonright W^\perp}$ i łączymy je do rozkładu całego $\rho$.
\end{proof}

Powstaje pytanie o unikalność rozkładu reprezentacji na reprezentacje nieprzywiedlnie. Wystarczy wziąć reprezentację trywialną, czyli wszystkie elementy grupy działają przez macierz identycznościową. Taka reprezentacja rozkłada się w sumę $1$-wymiarowych reprezentacji, a jak wiadomo na przestrzeni $n$-wymiarowej mamy niesłychanie wiele wyborów $n$ liniowo niezależnych jej podprzestrzeni $1$-wymiarowych.

\begin{theorem}{}{}
  Niech $G$ będzie grupą skończoną, a $\{(\rho_1, V_1),...,(\rho_k, V_k)\}$ wszystkimi jej nieizomorficznymi reprezentacjami nieprzywiedlnimi. Wtedy
  $$|G|=\sum (\dim V_i)^2$$
\end{theorem}

\begin{proof}
  Dowód z niezerowym prawdopodobieństwem pojawi się kiedyś. Opiera się on o fakt, że algebra grupowa jest izomorficzna z sumą prostą algebra macierzowych.
\end{proof}

\begin{definition}{reprezentacja regularna}{}
  % Reprezentacja regularna to działanie grupy na samej sobie przez lewe przesunięcia.
  Lewa regularna reprezentacja $\lambda$ grupy $G$ to reprezentacja na przestrzeni funkcji ciągłych, $C^*(G)$, zadana przez
  $$\lambda(g)f(x)=f(g^{-1}x)$$
\end{definition}

\begin{lemma}{}{}
  Niech $L:V\to W$ będzie dowolnym przekształceniem liniowym, a $\rho_1$, $\rho_2$ niech będą dwie reprezentacje $G$. Wówczas odwzorowanie $A_G(L):V\to W$ dane wzorem
  $$\frac{1}{|G|}\sum_{g\in G}\rho_1(g)^{-1}L\rho_2(g)$$
  jest przeplataczem $\rho_1$ i $\rho_2$.
\end{lemma}

\subsection{Produkt tensorowy reprezentacji}

Każdy zna i kocha produkt prosty dwóch przestrzeni wektorowych $V\oplus V'$, gdzie dodawanie i mnożenie przez skalary odbywa się po współrzędnych. 

\begin{definition}{produkt tensorowy}{}
  Produkt tensorowy przestrzeni $A$ i $B$ to jedyna przestrzeń wektorowa $A\otimes B$ taka, że każde dwuliniowe odwzorowanie $A\oplus B\to C$ faktoryzuje się przez $A\otimes B$. To znaczy komutuje diagram
  $$\begin{tikzcd}
    A\oplus B\arrow[r]\arrow[dr] & A\otimes B \arrow[d, dashed]\\ 
                                 & C
  \end{tikzcd}$$
  gdzie strzałka $A\oplus B\to A\otimes B$ to dwuliniowe "mnożenie" $(a,b)\mapsto a\otimes b$.
\end{definition}

Tak się składa, że jeśli $V$ ma bazę $\{v_i\}_{i\leq n}$, a $W$ - $\{w_i\}_{i\leq m}$, to ich iloczyn tensorowy $V\otimes W$ ma bazę $\{v_i\otimes w_j\}_{i\leq n, j\leq m}$.

\begin{example}[m]
  \item $\R\otimes \R\cong \R$ jako przestrzenie wektorowe nad $\R$: 
    $$(x,y)=(x\cdot 1, y\cdot 1)\mapsto (x\cdot 1)\otimes (y\cdot 1)=xy(1\otimes 1).$$
  \item Ogólniej: jeśli $V$ jest przestrzenią liniową nad ciałem $K$, to $V\otimes_KK\cong V$.
  \item Dla trzech przestrzeni liniowych $V$, $V'$, $W$ mamy $(V\oplus V')\otimes W\cong (V\otimes W)\oplus (V'\otimes W)$ - zupełnie jak mnożenie i dodawanie
  \item $\C\otimes_\R\C\cong \C^2$
\end{example}

Jeśli $\rho:G\to V$ i $\rho':G\to V'$ są reprezentacjami $G$, to możemy zdefiniować reprezentację (działanie) $G$ na $V\otimes V'$ w następujący sposób:
$$g(v\otimes v'):=gv\otimes gv'=\rho(g)v\otimes \rho'(g)v'.$$

Produkt tensorowy bardzo przypomina mnożenie liczb, ale nie do końca. Jeśli rozważamy $v\otimes w\in V\otimes V$, to bardzo rzadko $v\otimes w=w\otimes v$. Czasami jednak przemienność mnożenia wektorów byłaby przydatna - stąd definiujemy tensory symetryczne w którym relacja ta zachodzi zawsze.

\begin{definition}{produkt symetryczny}{}
  Dla przestrzeni wektorowej $V$ definiujemy jej produkt symetryczny jako
  $$S^2(V):=V\otimes V\langle v\otimes w-w\otimes v\rangle.$$
  To znaczy produkujemy odwzorowanie $V\otimes V\to V\otimes V$ posyłające $v\otimes w\to w\otimes v$ i wydzielamy przez jego jądro. Wtedy $v\otimes w=w\otimes v$.
  \bigskip
  
  Definicję tę można uogólnić: $n$-ty produkt symetryczny $V$ to
  $$S^n(V):=V^{\otimes n}/\langle v_1\otimes..\otimes v_i\otimes v_{i+1}\otimes...\otimes n_v-v_1\otimes...\otimes v_{i+1}\otimes v_i\otimes ... \otimes v_n\rangle.$$
\end{definition}

Tak jak istnieją macierze symetryczne i antysymetryczne, hermitowskie i antyhermitowskie, symetryczny produkt tensorowy też ma swojego kuzyna spod ciemnej gwiazdy, nazywającego się \buff{produktem alternującym}.

\begin{definition}{produkt alternujący}{}
  Dla przestrzeni wektorowej $V$ definiujemy jej produkt alternujący (potęga zewnętrzna) jako
  $$\Lambda^2(V):=V\otimes V\langle v\otimes w+w\otimes v\rangle.$$
  To znaczy produkujemy odwzorowanie $V\otimes V\to V\otimes V$ posyłające $v\otimes w\to -w\otimes v$ i wydzielamy przez jego jądro. Wtedy w ilorazie $v\otimes w=-w\otimes v$.
  \bigskip
  
  Definicję tę można uogólnić: $n$-ty produkt alternujący $V$ to
  $$\Lambda^n(V):=V^{\otimes n}/\langle v_1\otimes..\otimes v_i\otimes v_{i+1}\otimes...\otimes n_v+v_1\otimes...\otimes v_{i+1}\otimes v_i\otimes ... \otimes v_n\rangle.$$
\end{definition}

Jak się okazuje, symetria i antysymetria sumują się do całości, tzn.
$$V\otimes V\cong S^2(V)\oplus \Lambda^2(V)$$
przez operację symetryzacji i antysymetryzacji (?). Pierwsza z dowolnego tensora robi tensorek symetryczny
$$v\otimes w\mapsto\frac{1}{2}(v\otimes w+w\otimes v)$$
druga z dowolnego tensora robi tensor antysymetryczny
$$v\otimes w\mapsto\frac{1}{2}(v\otimes w-w\otimes v).$$
Sumując obie dostajemy $v\otimes w$ z powrotem.

\subsection{Charakter}

\begin{definition}{charakter}{}
  Niech $\rho:G\to GL(V)$ będzie reprezentacją grupy $G$ nad ciałem $K$. Funkcję $\chi:G\to K$ daną przez $\chi(g)=Tr(\rho(g))$ nazywamy \buff{charakterem reprezentacji $\rho$}.
\end{definition}

Innymi słowy: charakter przypisuje reprezentacji sumę jej wartości własnych.

Na tę chwilę będą interesować nas unitarne reprezentacje grup skończonych.

\begin{fact}{}{}
  Jeśli $\chi$ jest charakterem $n$-wymiarowej reprezentacji to
  \begin{itemize}
    \item $\chi(1)=n$
    \item $\chi(s^{-1})=\chi(s)^*$ (gdy reprezentacja jest unitarna)
    \item $\chi(tst^{-1})=\chi(s)$
  \end{itemize}
\end{fact}

\begin{lemma}{Schura}{}
  Niech $\rho_i:G\to GL(V_i)$, $i=1,2$ będą dwoma nieprzywiedlnimi reprezentacjami grupy $G$. Niech $f:V_1\to V_2$ będzie przeplataczem $\rho_1$ i $\rho_2$. Wtedy
  \begin{enumerate}
    \item jeśli $\rho_1$ nie jest izomorficzne z $\rho_2$ to $f=0$,
    \item jeśli $V_1=V_2$ oraz $\rho_1=\rho_2$ to $f$ jest wielokrotnością identyczności.
  \end{enumerate}
\end{lemma}

\begin{proof}
  \begin{enumerate}
    \item Rozważmy $x\in\ker f$, wtedy $0=\rho_2 fx=f\rho_1 x$, czyli $\rho_1x\in \ker f$, a więc $\ker f$ jest niezmiennicza na $\rho_1$. W takim razie $\ker f=V_1$ lub $\ker f=0$. Pierwszy przypadek jest przez nas pożądany, więc przyjrzyjmy się drugiemu. Jeśli $\ker f=0$, to $\im f$ jest niezmiennicze na działanie $\rho_2$ ($f\rho_1 x=\rho_2 fx$). W takim razie $\im f=V_2$ lub $\im f=0$. Skoro $\ker f=0$, to musowo $\im f=V_2$, a więc $f\neq 0$ musiałoby być izomorfizmem.
    \item Skoro pracujemy nad $\C$, to $f$ posiada co najmniej jedną wartość własną $\lambda$. Funkcja $f'=f-\lambda$ nie jest izomorfizmem, ale nadal jest przeplataczem $\rho$ z $\rho$, więc na mocy punktu pierwszego $f'=0$.
  \end{enumerate}
\end{proof}

Na przestrzeni zespolonych funkcji ze skończonej grupy $G$ możemy określić iloczyn skalarny wzorem
$$\langle \phi, \psi\rangle=\frac{1}{|G|}\sum_{g\in G}\phi(g)\psi(g)^*$$
W takim razie mając dwie unitarne reprezentacje $G$ jesteśmy w stanie
\begin{itemize}
  \item policzyć normę charakteru tej reprezentacji,
  \item sprawdzić, czy są one ortogonalne licząc iloczyn skalarny ich charakterów.
\end{itemize}

% Jak się okazuje, charaktery są przydatnym narzędziem do badania reprezentacji grup.


\begin{lemma}{}{}
  Reprezentacja $V$ o charakterze $\chi$ jest nieprzywiedlnia $\iff$ $\langle\chi,\chi\rangle=1$
\end{lemma}

\begin{proof}
  Niech $\rho$ będzie nieprzywiedlnią reprezentacją $G$ działającą przez macierze $(a_{ij})(g)$. Wówczas $\chi=\sum a_{ii}(g)$ oraz
  $$\langle \chi, \chi\rangle=\sum_{i,j} \langle a_{ii}(g),a_{jj}(g)\rangle=\sum_{i,j}a_{ii}(g)\overline{a_{jj}(g)}=\sum_{i,j}a_{ii}(g)a_{jj}(g^{-1})=\begin{cases}1 & i=j\\ 0 & i\neq j\end{cases}$$
\end{proof}

\begin{lemma}{}{}
  Nieprzywiedlnie nieizomorficzne reprezentacje mają ortogonalne charaktery.
\end{lemma}

\begin{proof}
  Niech $\rho_1$, $\rho_2$ będą dwiema nieprzywiedlnimi reprezentacjami, a $\chi_1$, $\chi_2$ ich charakterami.
\end{proof}

\begin{theorem}{}{}
  Niech $V=W_1\oplus W_2\oplus ... \oplus W_n$ będzie rozkładem reprezentacji $G$ o charakterze $\phi$ na nieprzywiedlnie podreprezentacje. Jeśli $W$ jest inną reprezentacją $G$ z charakterem $\chi$ to liczba $W_i$ które są z $W$ izomorficzne wynosi $\langle \chi, \phi\rangle$.
\end{theorem}



% \begin{theorem}{}{}
%   Nieredukowalne reprezentacje o tym samym charakterze są izomorficzne.
% \end{theorem}
%
% \begin{proof}
%   Jest to szybki wniosek z te
% \end{proof}



\section{Miara Haara}

Oryginalna praca w której Haar dowodzi istnienie lewo-niezmienniczej miary jest konstruktywna, ale my, ku smutkowi wykładowcy, podamy dowód bez liczenia.

\begin{definition}{}{}
  Niech $G$ będzie grupą topologiczną lokalnie zwartą. Miara Haara $\mu$ na $G$ to niezerowa miara Radona która jest lewo niezmiennicza, tj. $\mu(gE)=\mu(E)$ dla każdego $g\in G$ oraz dla każdego $E$ borelowskiego.
\end{definition}

Dla grup, które nie są lokalnie zwarte możemy definiować miary Gaussa.

\begin{lemma}{}{}
  Jeśli $\mu$ jest lewo-niezmienniczą miarą Haara na $G$ to $\tilde{\mu}(E)=\mu(E^{-1})$ jest prawo-niezmienniczą miara Haara na $G$.
\end{lemma}

\begin{proof}
  \begin{align*}
    \tilde{\mu}(Eg)=\mu(g^{-1}E^{-1})=\mu(E^{-1})=\tilde{\mu}(E)
  \end{align*}
\end{proof}

% Istnieje grupa, która jest $T_{5/2}$ ale nie $T_4$

\begin{theorem}{}{}
  Na lokalnie zwartej grupie istnieje jedyna z dokładnością do stałych lewo-niezmiennicza miara.
\end{theorem}

Pierwsze dowody istnienia miary Haara były dla grup ośrodkowych i my się do tych grup ograniczymy.

\begin{proof}
  Pominiemy dokładny dowód istnienia miary Haara. Idea polega na pokazaniu, że 
  \begin{itemize}
    \item przestrzeń funkcji nieujemnych o zwartym nośniku, $C_c^+(G)$, jest lokalnie wypukła,
    \item lokalnie zwarta grupa $G$ działa na niej przez operatory liniowe przez lewe przesunięcia, tzn. $(gf)(h)=f(g^{-1}h)$,
    \item istnieje zwarty zbiór $K\subseteq C_c^+(G)$ zachowywany przez działanie $G$
  \end{itemize}
  
  a następnie skorzystaniu z twierdzenia Kakutani, podanego niżej.

  \begin{theorem}{Kakutani}{}
    Niech $K\subseteq E$ będzie zwartym podzbiorem przestrzeni lokalnie wypukłej, czyli każdy otwarty podzbiór zawiera mniejszy otwarty podzbiór wypukły. Jeśli $G$ działa na $E$ przez operatory liniowe tak, że $gK\subseteq K$ dla wszystkich $g\in G$ to istnieje punkt $k_0\in K$ taki, że $gk_0=k_0$ dla wszystkich $g\in G$.
  \end{theorem}

  W ten sposób dostajemy funkcję $\mu$ będącą fix punktem działania $G$. Pozostaje pokazać (w ramach ćwiczenie dla czytelnika), że jest ona jedyna z dokładnością do skalowania.
\end{proof}


% na wykładzie - szkic dowodu dla grup zwartych
%
% on leci z riesza
%
% $A(S_3)$ - afiniczna grupa permutacji








